% ============================================================
%  Front matter: title, abstract, highlights, keywords, Section 1
% ============================================================

\begin{abstract}
Remaining useful life (RUL) predictions support maintenance decisions
only if their uncertainty is reported reliably, and most reliably in
the near-failure regime where decisions are taken. That regime is also
the least represented in run-to-failure data, and the few available
run-to-failure units leave little data for calibrating uncertainty.
This paper proposes a three-stage framework for bearing RUL prediction
that addresses these constraints together. A stage-conditioned
Wasserstein GAN synthesises complete multivariate near-failure windows
for training (C1). A multi-task BiLSTM with temporal attention predicts
the RUL and the degradation stage jointly from a shared encoder (C2).
Stage-conditional cross-conformal prediction (CV+) turns each point
prediction into an interval, uses every non-test bearing for
calibration, and conditions the interval on the stage predicted by
C2 (C3). Stage labels and RUL targets are derived automatically from a
guarded health-indicator criterion. On the XJTU-SY dataset, evaluated
with bearing-level cross-validation over five seeds, multi-task
learning lowers the near-failure RMSE from 0.225 to 0.206 and the
overall RMSE from 0.385 to 0.291. Cross-conformal calibration restores
nominal coverage: 0.917 overall at a 90\% target, against 0.74--0.83
for split conformal calibrated on two bearings. Conditioning on the
stage predicted by the multi-task model narrows the intervals of the
full framework by about 20\% relative to a single quantile, at higher
coverage.
Generative augmentation produces realistic windows but does not change
accuracy or interval quality significantly on its own. A per-bearing
analysis shows that coverage holds on average over bearings but not
for every bearing, with the shortest-lived bearings least well
covered.
\end{abstract}

\begin{highlights}
\item Stage labels and piecewise RUL targets derived without manual annotation.
\item Multi-task RUL and stage learning lowers near-failure and overall RMSE.
\item Cross-conformal calibration uses all non-test bearings and restores 90\% coverage.
\item The predicted stage narrows conformal intervals by about 20\% at equal coverage.
\item Per-bearing coverage exposes the limits of the guarantee on short-lived bearings.
\end{highlights}

\begin{keyword}
Remaining useful life \sep Rolling bearings \sep Conformal prediction
\sep Uncertainty quantification \sep Multi-task learning
\sep Generative adversarial networks
\end{keyword}

\end{frontmatter}

% ------------------------------------------------------------
\section{Introduction}
\label{sec:introduction}

Rolling bearings are among the most frequent sources of failure in
rotating machinery, and their unplanned failure stops production and
can damage surrounding components. Predictive maintenance addresses
this by estimating the remaining useful life (RUL) of each bearing from
its condition-monitoring signals and scheduling intervention before
failure. Deep learning models trained on vibration data now estimate
bearing RUL with good average accuracy on public run-to-failure
benchmarks such as XJTU-SY~\cite{8576668, lei2019xjtusy}. A point
estimate alone, however, is a weak basis for a maintenance decision.
The decision maker must also know how far the true RUL may lie from the
estimate, and this knowledge matters most close to failure, when the
cost of an error is highest.\\

Three properties of run-to-failure data make this difficult. First,
bearings spend most of their life in a healthy state, so the
near-failure regime, where predictions matter most, is the least
represented in the training data. Second, most uncertainty estimates
used for RUL, such as Bayesian layers, Monte Carlo dropout or Gaussian
output layers, rest on distributional assumptions and carry no
coverage guarantee; they are frequently overconfident in
practice~\cite{wang_uncertainty-aware_2026}. Conformal
prediction~\cite{vovk2005algorithmic, lei2018distribution} removes
these assumptions and provides intervals with finite-sample coverage,
but its guarantee rests on the calibration data resembling the test
data. Third, a run-to-failure dataset contains few units: XJTU-SY has
fifteen bearings. Holding out a calibration set therefore leaves very
few bearings to represent the variability between units, and the
coverage guarantee, which is formally marginal, may not hold where it
is needed.\\

Each of these difficulties has been addressed separately. Generative
adversarial networks have been used to enlarge the near-failure
regime~\cite{9486650, lu_joint_2022}, mostly on scalar health
indicators and evaluated through point accuracy only. Multi-task
models that predict the degradation stage alongside the RUL improve the
learned representation~\cite{8709822}. Conformal prediction has very
recently been applied to bearing RUL~\cite{wang_uncertainty-aware_2026,
11373138}, with coverage reported in aggregate over all windows. No
existing study combines the three, and none conditions the conformal
interval on the degradation stage that a multi-task model predicts, or
addresses the small number of units available for calibration.\\

This paper proposes a three-stage framework that addresses these
constraints together, and makes the following contributions:
\begin{itemize}
  \item \textbf{C1, stage-conditioned generative augmentation.} A
  Wasserstein GAN with gradient penalty, conditioned on the degradation
  stage, synthesises complete multivariate near-failure windows, which
  enter the training data only.
  \item \textbf{C2, multi-task RUL and stage prediction.} A BiLSTM with
  temporal attention predicts the RUL and the degradation stage from a
  shared encoder. The stage labels are derived automatically from a
  guarded health-indicator criterion, and the RUL target is piecewise.
  \item \textbf{C3, stage-conditional cross-conformal prediction.}
  Cross-conformal calibration (CV+) uses every non-test bearing for
  calibration, and the interval is conditioned on the stage predicted
  by the multi-task model of C2, so that a single network supplies the
  RUL estimate, the stage and the interval.
\end{itemize}
The framework is evaluated on XJTU-SY with bearing-level
cross-validation, a $2 \times 2$ ablation of C1 and C2, five random
seeds, and a per-bearing analysis of coverage. The evaluation is
reported with its full protocol (RUL scale, target shape, scoring
window and aggregation), so that the results can be compared with the
literature without ambiguity.\\

The remainder of the paper is organised as follows.
Section~\ref{sec:related_work} reviews related work.
Section~\ref{sec:methodology} presents the framework.
Section~\ref{sec:experimental_setup} describes the dataset, the data
preparation and the evaluation protocol. Section~\ref{sec:results}
reports the results, Section~\ref{sec:discussion} discusses them and
their limitations, and Section~\ref{sec:conclusion} concludes.
